\chapter{Distribución normal multivariante en alta dimensión}

En este capítulo, el objetivo será completar las pruebas realizadas en la Sección 2.1 para los teoremas 2.2, 2.4 y 2.6 dando resultados que nos aseguren que los valores obtenidos para las esperanzas de $\|X\|$, $\|X - Y\|$, e $\langle X, Y \rangle$, si $X, Y \sim \mathcal{N}(0_d, I_d)$ e independientes, son también aplicables a muestras aleatorias de estas distribuciones.

\begin{lemma}[Desigualdad de Bernstein]
Sea $(\Omega, \sigma, P)$ un espacio probabilístico y $X_1, \dots, X_d : \Omega \to \mathbb{R}$ variables aleatorias independientes dos a dos con $E(X_l) = 0$ y momentos de orden $k$ acotados por $|E(X_l^k)| \leq \frac{k!}{2}$ para todo $l = 1, \dots, d$ y $k \geq 2$. En estas condiciones, para todo $\epsilon > 0$, se tiene que:
\[ P [ |X_1 + \dots + X_d| \geq \epsilon ] \leq 2e^{-\frac{1}{4} \min(\frac{\epsilon^2}{d}, \epsilon)}. \]
\end{lemma}

\begin{theorem}[Teorema del Anillo Gaussiano]
Sean $(\Omega, \sigma, P)$ un espacio probabilístico, $X : \Omega \to \mathbb{R}^d$ una variable aleatoria con distribución $\mathcal{N}(0_d, I_d)$ y $0 \leq \epsilon \leq \sqrt{d}$. Entonces se cumple que
\[ P [ \|X\| - \sqrt{d} \leq \epsilon ] \geq 1 - 2e^{-\epsilon^2/16}. \]
\end{theorem}

(El resto del capítulo 4 se puede añadir de forma similar)
